\ChapterImageStar[cap:caracterizacion-universos]{Caracterización de los Universos HTCondor}{./images/fondo.png}\label{cap:caracterizacion-universos}
\mbox{}\\

En esta sección se presenta una descripción detallada de la herramienta ArchiMate, considerada un referente conceptual y metodológico fundamental para el desarrollo de la Notación de Diseño de Infraestructura de~\TI~(\NDITI). A través de esta caracterización se analizan sus principales elementos, capas, relaciones y principios de modelado, así como su potencial de adaptación a los requerimientos específicos de NDITI. De igual manera, se revisan sus fortalezas y limitaciones, con el fin de determinar en qué medida esta herramienta puede servir como base metodológica para la definición y representación visual de infraestructuras complejas en entornos académicos o institucionales.


\section{ArchiMate}

Como se ha mencionado anteriormente, ArchiMate es un lenguaje de modelado abierto e independiente, gestionado como estándar técnico por The Open Group, diseñado específicamente para la descripción, el análisis y la visualización de arquitecturas empresariales~\cite{ArchiMate}. Uno de sus propósitos principales es la definición de las relaciones entre conceptos pertenecientes a diferentes dominios arquitectonicos, estos se ubican en un punto intermedio entre los conceptos ams detallados usados para modelar dominios individuales~\cite{Stojanovic01}. Este lenguaje ha estado en constante evolución, cuya última versión ``ArchiMate 3.2'' fue publicada en Octubre del 2022~\cite{ArchiMate}.

En el marco de la NDITI, ArchiMate constituye un referente metodológico y formal para la especificación visual de diseños de infraestructura, ya que provee un lenguaje que facilita la interoperabilidad y el entendimiento común entre diseñadores, ingenieros y gestores de TI.


\subsection{\textit{Objetivos Principales}}

\begin{itemize}
	\item Proveer un lenguaje común, independiente de herramientas y estandarizado para describir la arquitectura de una organización, facilitando que sus diferentes dominios se modelen de forma coherente~\cite{ArchiMate}.
	\item Facilitar la consistencia y coherencia en el modelado de arquitectura a traves de la alineación entre negocio, aplicaciones y tecnología mediante una representación clara concetrada en un solo framework~\cite{OMG01}.
	\item Facilitar la toma de decisiones informadas mediante el análisis de decisiones arquitectónicas, evaluación de riesgos y evaluación de escenarios alternos~\cite{archimetric01}.
	\item Ofrecer un marco visual estandarizado e inequívoco de modo que los stakeholders de perfiles, tanto técnicos como no técnicos, puedan colaborar y comunicarse entre si garantizando la comprensión y la alineación compartidas~\cite{archimetric01}.
\end{itemize}


\subsection{Estructura del Lenguaje}

Segun el documento de especificación de~\cite{ArchiMate}, este mismo define estructuras de elementos y realciones agrupadas y especializadas en diferentes capas específicas, las cuales seran detaladas a continuación:

\begin{itemize}
	\item \textbf{Capa de Negocio:} Referente a los servicios ofrecidos a los consumidores de una organización, los cuales a su vez son realizados por los procesos de neogocio al interior de la misma.
	\item \textbf{Capa de Aplicación:} Describe los servicios de aplicación que dan soporte a los procesos de negocio de la organización y a su vez las aplicaciones que los implementan.
	\item \textbf{Capa de tecnología:} Esta capa agrupa tanto las tecnologías de operación, como las pertenecientes a \TI. Permitiendo incluir tanto la tecnología que da soporte a las aplicaciones y capas de negocio, como también instalaciones y equipos físicos.
\end{itemize}

Sin embargo, Archimate tambien cuenta con un framework extendido llamado full framework, incorporando nuevas capas y perspectivas adicionales buscando un efoque mas integral en el modelado~\cite{leanix01}. EL documento de especificacion mencionado anteriormente (\cite{ArchiMate}) también hace énfasis en las capas adicionales:

\begin{itemize}
	\item \textbf{Capa de Estrategia:} Agrupa los elementos usados para modelar la direccion estrategica de la organización en lo que respecta a su impacto en la arquitectura. Su diferencia principal con la capa de negocio radica en que esta última modela únicamente la organización operativa de la empresa.
	\item \textbf{Capa de implementación/Migración:} Soporta todo aquello relacionado con la implementación y migración de arquitecturas, incluyendo programas y proyectos de respaldo para diferentes aspectos organizativos (Programas, portafolios y proyectos) y su planificación.
\end{itemize}

Tomando en cuenta ambas estructuras por capas mencionadas anteriormente, la jerarquía del modelo full framework al completo es la siguiente:
 
\begin{itemize}
	\item \textbf{Capa de Estrategia}
	\item \textbf{Capa de Negocio}
	\item \item \textbf{Capa de Aplicación}
	\item \item \textbf{Capa de tecnología}
	\item \textbf{Capa de implementación/Migración}
\end{itemize}


\subsection{\textit{Elementos y relaciones clave}}

En su docuemnto de especificación,~\cite{ArchiMate}~define también un conjunto limitado y coherente de elementos básicos y de relaciones que permiten modelar la arquitectura de la empresa de forma sistemática y comprensible. Estos elementos se presentan se presentan en la tabla~\ref{table:archimate-elements}.

	\subsubsection{Elementos básicos por dominio}

	\begin{table}[H]
	\centering
	\renewcommand{\arraystretch}{1.3}
	\setlength{\tabcolsep}{8pt}
	\begin{tabular}{@{}p{4cm}p{6cm}p{6cm}@{}}
	\toprule
	\textbf{Capa / Aspecto} & \textbf{Elementos principales} & \textbf{Relaciones típicas} \\ 
	\midrule
	\textbf{Estrategia / Motivación} &
	\textit{Driver}, \textit{Goal}, \textit{Outcome}, \textit{Principle}, \textit{Requirement}, \textit{Capability}, \textit{Resource} &
	\textit{Influence}, \textit{Realization}, \textit{Aggregation} \\ 
	\midrule
	\textbf{Negocio} &
	\textit{Business Actor}, \textit{Business Role}, \textit{Business Process}, \textit{Business Service}, \textit{Business Object} &
	\textit{Assignment}, \textit{Serving}, \textit{Flow}, \textit{Composition} \\ 
	\midrule
	\textbf{Aplicación} &
	\textit{Application Component}, \textit{Application Interface}, \textit{Application Service}, \textit{Data Object} &
	\textit{Realization}, \textit{Serving}, \textit{Access}, \textit{Flow} \\ 
	\midrule
	\textbf{Tecnología / Infraestructura} &
	\textit{Node}, \textit{Device}, \textit{System Software}, \textit{Artifact}, \textit{Technology Service} &
	\textit{Composition}, \textit{Assignment}, \textit{Realization}, \textit{Flow} \\ 
	\midrule
	\textbf{Implementación / Migración} &
	\textit{Deliverable}, \textit{Work Package}, \textit{Plateau}, \textit{Gap} &
	\textit{Triggering}, \textit{Flow}, \textit{Aggregation} \\ 
	\bottomrule
	\end{tabular}
	\caption{Elementos y relaciones clave del lenguaje ArchiMate según sus diferentes capas.}
	\label{table:archimate-elements}
	\end{table}


\section{Limitaciones de ArchiMate}
Si bien ArchiMate se ha consolidado como uno de los lenguajes de modelado más utilizados para representar arquitecturas empresariales de forma integrada y estandarizada, Diversos estudios y experiencias en la industria han identificado limitaciones relacionadas con la complejidad de su aprendizaje, el nivel de detalle técnico que ofrece, su enfoque predominantemente organizacional y las condiciones de uso de las herramientas que lo implementan, como se presenta a continuación.

\begin{itemize}
	\item \textbf{Curva de aprendizaje elevada:} Una revisión realizada por el grupo BOC, indica que ArchiMate es un lenguaje de modelado complejo, debido a que el mismo incluye numerosas capas, elementos y relaciones~\cite{BOC01}. Lo anterior exige que los stakeholders y modeladores tengan una comprensión sólida de ArchiMate y de la arquitectura empresarial para poder entender estos modelos y comunicarse de manera efectiva.
	\item \textbf{Menor precisión técnica:} Análisis realizados por la empresa de consultoria NILUS señalan que si bien ArchiMate es muy adecuado para modelar a nivel arquitectónico, el mismo presenta limitaciones para describir detalles de infraestructura técnica debido a su carencia de estructuras sólidas para la arquitectura de datos y capacidades formales de modelado entidad-relación~\cite{Nilus01}.
	\item \textbf{Enfoque organizacional:} La meta principal de ArchiMate es ser lo mas pequeño posible~\cite{ArchiMate}. Por lo que está diseñado principalmente para describir la arquitectura empresarial y no para escenarios exclusivos de infraestructura donde puede no tener las herramientas necesarias para la comunicación efectiva de estos escenarios.
	\item \textbf{Dependencia de herramientas propietarias:} Aunque el lenguaje es abierto, múltiples estudios y análisis del mercado señalan que las herramientas de arquitectura empresarial implican costes elevados, de implementación y mantenimiento cual puede hacer difícil su justificación en organizaciones pequeñas o con baja madurez en arquitectura\cite{kenresearch01}.
	\item \textbf{Mantenimiento y escalabilidad:} En la revisión del grupo BOC anteriormente mencionada, tambien se expone la dificultad de mantener la consistencia y la precisión que aumenta cuando los modelos crecen en tamaño y complejidad, lo que en ultima instancia genera modelos difíciles de gobernar o interpretar al aumentar el riegos de errores e inconsistencias~\cite{BOC01}.
	\item \textbf{Ambigüedad semántica:} En algunos estudios se ha observado que ArchiMate puede ser multi-interpretable, con conceptos importante ausentes y sin ofrecer metodos estructurados para capturar reglas de negocio~\cite{Krouwel01}. Lo que puede restar eficacia en la comunicación o en la toma de decisiones basada en el modelo.
\end{itemize}

Visto lo anterior, se puede concluir que ArchiMate constituye una herramienta valiosa para la representación y modelado de infraestructuras, presenta limitaciones que deben considerarse en su aplicación práctica. Su curva de aprendizaje elevada y la escasa precisión en aspectos técnicos de bajo nivel pueden dificultar su adopción. Además, su cobertura conceptual no abarca todas las perspectivas posibles —como las operativas, de seguridad o de configuración detallada—, lo que limita su utilidad cuando se requiere una representación exhaustiva de la infraestructura o de dominios técnicos específicos.

\subsection{Aplicaciones prácticas y herramientas compatibles}

El uso de ArchiMate no se limita a sus especificaciones teóricas: existen múltiples herramientas (libres y comerciales) que permiten modelar conforme a su metamodelo y aplicar la notación en contextos reales. A continuación se presenta una selección de herramientas relevantes:

\begin{itemize}
	\item \textbf{Archi:} Conjunto de herramientas de modelado enfocada en la accesibilidad. Soporta la creación de modelos ArchiMate completos permitiendo ser operada por modeladores y arquitectos empresariales de todos los niveles. Su naturaleza es de código abierto y multiplataforma~\cite{Archi01}.
	\item \textbf{BiZZdesign Enterprise Studio:} Se presenta como una solución corporativa avanzada para arquitectura empresarial, su soporte del lenguaje ArchiMate 3.0 es oficialmente certificado por The Open Group, permitiendo herramientas de colaboración, dashboards, reporting y modelado integrado~\cite{Bizz01}.
	\item \item \textbf{Sparx Enterprise Architect:} Herramienta multipropósito con soporte para diversos lenguajes de modelado, entre ellos ArchiMate cuyo soporte es oficial. Al ser enfocada en multiples lenguajes permite operaciones como modelado completo, trazabilidad entre capas, vistas y almacenamiento en repositorio~\cite{EA_GUIDE01}.
\end{itemize}

Las herramientas mencionadas anteriormente estan enfocadas en facilitar la adopción de ArchiMate en organizaciones de dsitintos niveles de tamaño o madurez, permitiendo construir repositorios de arquitectura, definir vistas para stakeholders, realizar análisis de impacto y generar artefactos de modelado para las mismas.

Una vez terminada la caracterización de la herramienta ArchiMate, se puede evidenciar que, a pesar de ofrecer una estructura sólida para modelar arquitecturas empresariales, presenta limitaciones al representar dominios técnicos o de infraestructura de TI, donde puede resultar ambiguo o insuficientemente detallado. Su curva de aprendizaje elevada, la falta de precisión técnica y cierta ambigüedad semántica reducen su eficacia en escenarios que requieren especificaciones más formales. Estas limitaciones evidencian la necesidad de una notación complementaria, como la NDITI, que permita una representación más clara y ajustada al contexto de diseño de infraestructuras tecnológicas.